\section{Parameter choice and completion of the local estimate}\label{sec:completion}

The preceding sections provide the coefficient-dependent inputs to
the completion arguments in \cite[Section~10]{H} and
\cite[Section~7]{V}. We now explain the common parameter choice and
the two distinct conclusions obtained from it. First we bound the
fixed-input constant $C_N(q)$ independently of $N$. Only after this
absorption do we complete the proof for depth-dependent inputs.

\subsection{The retained-cell estimate}
We use the base reductions of \cite[Proposition~5.8]{H} and
\cite[Proposition~5.6]{V}, with the modifications proved above.
At one noninitial accuracy level their target, on a normalized root
$R$, is a sum of terms
$\abs I\min\{\eta,\abs{H_I(z_I)}\}$ over a specified subset of
descendant cells. Here $z_I$ is the four-tuple supplied by the base
expansion at $I$, and $\eta=\exp(-k^\theta)$ is its quality parameter,
with $k$ the accuracy log and $0<\theta<1$ fixed. The constructions
then decompose each of these four-tuples into the localized terms
estimated in the preceding sections.

The stopping constructions in \cite[Section~10.1]{H} and
\cite[Section~7.1]{V} are used in their stated order. Main columns and
their uses are fixed first; mode expansions and the whole side
functions are fixed next; helper projections are constructed after
those side functions. Size and counting stops are then imposed on
these fixed objects. A stop removes original summands from an
attempt; it does not redefine a previously fixed stack function.
In particular, the maximal functions in the short-block estimates
remain defined on whole retained cells, including the portions below
later stopping cells.

Let $\mathcal A_R^{\mathrm{ret}}$ denote the cells not removed by
these stops in the attempt rooted at $R$. The conclusions of the
preceding sections give the same estimate as the retained-cell
calculations in the two companions:
\begin{equation}\label{eq:retained}
\begin{aligned}
 \sum_{I\in\mathcal A_R^{\mathrm{ret}}}\abs I
       \min\{\eta,\abs{H_I(z_I)}\}
 &\le \Bigl[C_q\eta^{\varepsilon}P^{C_4}
       +(1+C_N(q))C_qP^{-1/2}G_0
                    (K+2)^{C_{19}(q-2)}\Bigr]\abs R\\
 &\quad+O(e^{-k}\abs R).
\end{aligned}
\end{equation}
Here $\varepsilon>0$ is fixed, $P$ is the transfer parameter, $K$
is the depth-block length, and $G_0$ is the row and stack overhead.
The constants may depend on $b,\delta_0,C_0$.

For clarity, the estimate has three sources. Terms with at least
three core factors use the powered core estimate, and terms with at
least three side factors use their square-energy bounds. With
$p=1-\varepsilon>2/3$, both are converted to the first term
of~\eqref{eq:retained} by
$\min\{\eta,a\}\le\eta^{1-p}a^p$.
Terms with exactly two sides use the transfer estimate for the
whole complementary arguments. Its main term is again included in
the first term. The remaining coupled near-band sum is estimated
at power one by the short-block comparison; it gives the second
term, including its factor $1+C_N(q)$, by~\eqref{eq:near-band-b}.
Physical and numerical
errors are estimated at power one as well. These alternatives
exhaust the four-position expansion. The minimum is applied to
the whole contribution of a cell before subadditivity is used;
no smallness of a separately selected group is being assumed.

\subsection{Restarts and the complexity hierarchy}
The integrated size, path-count, and square-energy bounds allow the
stopping thresholds to be enlarged by fixed factors so that the first
stopped cells occupy at most $\abs R/4$. At each such cell send the
original summands there and below to a new attempt. The local
normalizations and column bounds are available on every new root.
Linear-side energy is proved anew by the construction on that root,
not by asserting a proportional restriction of a global energy bound.
The resulting attempt roots satisfy
\begin{equation}\label{eq:attempt-packing}
 \sum_{R\text{ an attempt root below }R_0}\abs R
 \le\sum_{m=0}^{\infty}4^{-m}\abs{R_0}
 =\frac43\abs{R_0}.
\end{equation}
All the trees are finite, so the restart argument can equivalently be
read as induction on their remaining depths.

In the scheduled-input construction, the original depths, functions,
and block partitions remain attached to the summands. Restriction to
a new root preserves the pointwise size and variation bounds with
the same $W_j$. The row estimates of \cite[Section~4]{V} therefore
apply at each restart with no factor depending on $W_j$.

The distinction between the two complexity scales is essential.
The bounds supplied by the constructions are
\begin{equation}\label{eq:budget-hierarchy}
\begin{aligned}
 \log D_{\mathrm m}&\le C_q k^{C},
 &\log D_*&\le C_q(1+k+P)^C,\\
 K&\le C_q(1+k+P)^{C_5},
 &G_0&\le\exp\bigl(C_q(1+\log(2k))^C\bigr).
\end{aligned}
\end{equation}
Here $D_{\mathrm m}$ bounds the main and helper catalogs, and
$D_*$ bounds the larger mode, Fourier, and rational-height data.
The latter data enter the counting and mode estimates, but are not
fed back into the main catalog construction. Strong stack estimates
compress main projection families only; mode projections have
already been included in the Bessel rows. This is precisely the
separation in \cite[Section~10.3]{H} and \cite[Section~7.3]{V}.

The coefficient modifications preserve this separation. Chart
substitutions multiply heights and derivative bounds by fixed
slope-dependent factors; the change of height base occurs only
in the mode construction. The sparse-count estimate retains its
dimension-polynomial reserves and its power saving. Thus all the
degrees $C,C_4,C_5,C_{19}$ can be fixed \emph{before} choosing $q$.
Their numerical values may be enlarged to cover both the fixed-input
and scheduled-input constructions. Only coefficients such as $C_q$
are allowed to depend on the later choice of $q$.

Choose $q\in(2,3)$ sufficiently close to $2$ that
\begin{equation}\label{eq:common-q}
 C_{19}C_5(q-2)<\frac18
\end{equation}
for both constructions. This common choice is legitimate because
local detection is available for every $q>2$ and short-block
comparison for every $2<q<3$. Next put
$P=\lceil\eta^{-\alpha}\rceil$, choosing a fixed $\alpha>0$ so small
that $\alpha C_4<\varepsilon/2$, as well as the analogous finitely
many upper bounds already included in the powered costs.
Finally choose the starting accuracy sufficiently large. Since $P$
dominates every fixed power of $k$ and $G_0=P^{o(1)}$, we obtain
\[
 P^{-1/2}G_0(K+2)^{C_{19}(q-2)}\le C_qP^{-1/4},\qquad
 \eta^{\varepsilon}P^{C_4}\le C_q\eta^{\varepsilon/2}.
\]
Also $P^{-1/4}\le\eta^{\alpha/4}$, and $e^{-k}$ is smaller than
these quality powers for large $k$. Consequently, for some $c_*>0$,
\eqref{eq:retained} and~\eqref{eq:attempt-packing} give the full
base-segment bound
\begin{equation}\label{eq:quality-gain}
 \sum_{I\text{ in the base segment}}\abs I
       \min\{\eta,\abs{H_I(z_I)}\}
 \le C_q\eta^{c_*}(1+C_N(q))\abs{R_0}.
\end{equation}

\subsection{Fixed inputs: accuracy summation and absorption}
The base accuracy levels satisfy $k_{t+1}=k_t^A$ for a fixed $A>1$.
For tuples whose later level is $t$, there are $O(t)$ choices of
the earlier level and a bounded number of slot patterns.
The segment-root packing in \cite[Section~5]{H}, followed by
\eqref{eq:attempt-packing}, gives a bounded total root length for
each tuple. The noninitial terms in~\eqref{eq:quality-gain} therefore
have total coefficient controlled by
\begin{equation}\label{eq:accuracy-series}
 C_q\sum_{t\ge12}t\exp(-c_*k_t^\theta).
\end{equation}
This series is finite, and its value is arbitrarily small when the
starting accuracy is sufficiently large. The replacement and
two-increment errors have the same summable multiplicity.
The finitely many initial levels have depth-independent bounds from
the active-scale reduction, as in \cite[Section~10.4]{H}.
Thus the modified construction proves the following normalized
conclusion.

\begin{proposition}[Normalized absorption]\label{prop:absorption}
For the fixed $q$ chosen above and every sufficiently small
$\tau>0$, there is $A_\tau<\infty$, independent of $N$, such that
the following holds. Suppose a finite rooted dyadic subtree has root
$I_*$, its summation cells lie in at most $N$ consecutive depths,
and its fixed scalar inputs satisfy
\[
 \left(\abs I^{-1}\int_I\abs{f_j}^q\right)^{1/q}\le1
 \quad(0\le j\le3, I\text{ in the subtree}).
\]
Then, for any subset of summation cells and admissible kernels,
\[
 \sum_{I\text{ a summation cell}}\abs I\,\abs{H_I(f_0,f_1,f_2,f_3)}
 \le (A_\tau+\tau C_N(q))\abs{I_*}.
\]
\end{proposition}

The argument just given is the proof: take the coefficient
in~\eqref{eq:accuracy-series} below $\tau$ and include its constant
part, the initial levels, and the numerical errors in $A_\tau$.
The base expansion is first taken at a finite accuracy cutoff.
At each fixed cell the terminal residual tends to zero by the
pair-comparison estimates of the base reduction; since there are
only finitely many cells, the passage to all accuracy levels is
legitimate. This is the concluding expansion step in
\cite[Section~10.4]{H}.

We include the outer normalization argument to make the absorption
explicit. Start with disjoint roots at the largest scale in a finite
collection from~\eqref{eq:CN}. At a root $J$, write
$a_{j,J}=(\abs J^{-1}\int_J\abs{f_j}^q)^{1/q}$. If some $a_{j,J}$
vanishes, all forms below $J$ vanish. Otherwise stop at the first
strict descendants $J'$ on which an average of $\abs{f_j}^q$
exceeds $16a_{j,J}^q$. The union of first stops has measure at most
$\abs J/4$. Dividing $f_j$ by $16^{1/q}a_{j,J}$ normalizes the
intervening tree, so Proposition~\ref{prop:absorption} bounds its sum
by
\[
 16^{4/q}(A_\tau+\tau C_N(q))\abs J\prod_ja_{j,J}.
\]
Repeat at the stopped roots. The sets $E_J$ left after deleting their
first stops are disjoint, have measure at least $3\abs J/4$, and
satisfy $M_qf_j\ge a_{j,J}$ on $E_J$. Hence
\[
 \sum_J\abs J\prod_ja_{j,J}
 \le\frac43\int_\R\prod_jM_qf_j.
\]
Taking the supremum in the definition of $C_N(q)$ gives
\[
 C_N(q)\le C A_\tau+C\tau C_N(q),
\]
where $C$ is independent of $N$ and bounded for $2<q<3$.
Choose $C\tau<1/2$ and use the already established finiteness of
$C_N(q)$ to absorb the last term. We have proved
\begin{equation}\label{eq:uniform-CN}
 C_{\mathrm H}(q):=\sup_{N\ge1}C_N(q)<\infty.
\end{equation}

\subsection{Depth-dependent inputs}
Now perform the scheduled-input construction with the same $q$,
using its own subsequent choices of $\alpha$ and starting accuracy.
Insert~\eqref{eq:uniform-CN} into~\eqref{eq:quality-gain}. The
factor $1+C_N(q)$ is now a fixed constant. The scheduled base
expansion of \cite[Section~5]{V} has the same accuracy multiplicity
and segment-root packing, with the depth-row estimates supplying
the uniform bounds. Its noninitial terms therefore sum
by~\eqref{eq:accuracy-series}. Its finitely many initial levels
are bounded by its active-scale reduction independently of the
number of depths and of every $W_j$.

Finally, at a fixed original cell, all accuracy updates are evaluated
on its same actual scheduled input $z_{j,s(I)}$. The per-cell pair
comparisons therefore make the terminal residual tend to zero.
There are finitely many cells, so the finite-cell sum commutes with
this limit. The triangle inequality and the uniform summed bounds
give~\eqref{eq:local-estimate}, exactly as in
\cite[Section~7.4]{V}. This proves Theorem~\ref{thm:local}, and hence,
by Section~\ref{sec:transfer}, Theorem~\ref{thm:main}.
